\section{Uživatelské požadavky}
Následující podkapitola se zabývá analýzou funkčních a nefunkčních požadavků kladených na novou aplikaci. Požadavky primárně vycházejí z rozhovoru se zástupcem firmy Ackee, která hledá způsob, jak měřit a porovnávat výkonnost vývojářů.

\subsection{Funkční požadavky}
Funkční požadavky popisují požadované funkcionality systému. Mohou například specifikovat, jakým způsobem bude uživatel moci pracovat se systémem, jaké procesy by měl systém podporovat a jaké budou vstupy a výstupy těchto procesů.

\begin{enumerate}[label=\textcolor{decoration}{\textbf{F\arabic*}}, leftmargin=9mm]
    \myItem{definice a správa sledovaných uživatelů}
    Systém bude umožňovat párování uživatelů s uživateli z platforem GitHub a GitLab. K jednomu uživateli může být ve vytvářené aplikaci přiřazeno více referencí na uživatele z platforem GitHub a GitLab. Uživatele též bude možné upravovat a mazat.

	\myItem{získání dat o uživatelích}
    V systému bude možné nastavit z jakých projektů a za jaké časové období se mají získávat data o sledovaných uživatelích. Data se následně budou automaticky stahovat ze systémů GitHub a GitLab a bude je možné získávat i zpětně (tj. i za období předcházející definici uživatelů a projektů ve vytvářené aplikaci).

    \myItem{definice časově omezených výzev}
    Administrátor systému bude moci definovat krátkodobé a dlouhodobé výzvy, za které uživatelé budou moci získávat „odznáčky“.

	\myItem{grafická prezentace dat}
    Data získaná ze systémů GitHub a GitLab budou ve webové aplikaci přehledně prezentována ve formě grafů nebo tabulek. Tyto přehledy budou dostupné jak pro jednotlivé uživatele, tak i pro definované skupiny uživatelů. Přehledy bude možné filtrovat dle období, projektu a platformy.
\end{enumerate}

\subsection{Nefunkční požadavky}
Nefunkční požadavky se na rozdíl od funkčních požadavků zabývají vlastnostmi a omezeními daného systému. Mohou se mezi nimi objevit požadavky, které budou podstatné jak pro samotné uživatele systému, jako například doba odezvy, zabezpečení systému, způsob zobrazování na mobilních zařízeních, tak zde mohou být uvedeny i požadavky techničtějšího charakteru, jako například použití konkrétních technologií.

\begin{enumerate}[label=\textcolor{decoration}{\textbf{N\arabic*}}, leftmargin=9mm]
    \myItem{responzivita}
    Webová aplikace bude přizpůsobovat svůj vzhled na základě rozlišení klientského zařízení.

    \myItem{zabezpečení}
    Komunikace s aplikací bude probíhat přes šifrovaný protokol HTTPS, hesla uživatelů budou ukládána pomocí hashů (otisků), budou implementovány ochrany proti XSS, SQL injection, včetně dalších bezpečnostních mechanismů.

    \myItem{výkon a dostupnost}
    Aplikace bude schopná v jednu chvíli obsluhovat jednotky až nižší desítky uživatelů s odezvou v řádu sekund.

    \myItem{právní požadavky}
    Sledovaného nebo registrovaného uživatele bude možné anonymizovat v souladu s obecným nařízením o ochraně osobních údajů (GDPR) nebo ho bude možné odstranit včetně všech jeho osobních dat.
\end{enumerate}
